\documentclass[11pt]{article}

\usepackage[margin=1in]{geometry}
\usepackage{amsmath,amssymb}
\usepackage{xcolor}
\usepackage{graphicx}
\usepackage{hyperref}
\usepackage{booktabs}
\usepackage{float}
\usepackage{url}

\definecolor{accent}{HTML}{2E75B6}

\title{Reward Shaping Ablation Study\\
\large From Local Penalty to Episode-Level Savings}
\author{Pulak Mehrotra}
\date{April 2026}

\begin{document}
\maketitle
\tableofcontents
\newpage

\bigskip\noindent
This document defines six reward formulations as a structured ablation
study, progressing from the narrowest possible signal (single worst MPD)
to the true episode-level objective augmented with optimal-gap shaping.
Each level expands the scope of information the agent receives, and the
ablation is designed to reveal which level is necessary and sufficient to
beat greedy.

Sections~\ref{sec:reachability}--\ref{sec:secondary} establish the three
conceptual building blocks: the topology constraints that limit action scope
(MPD reachability), the lower bound that anchors our shaping (optimal gap),
and the literature-derived techniques that improve any base reward (PBRS,
sub-episodes, counterfactual baselines).
Section~\ref{sec:ablation} presents the six concrete formulations.

% =========================================================================
\section{Background and Motivation}
\label{sec:background}

The agent distributes each arriving VM's memory request across
CXL-connected Memory Pool Devices (MPDs) in a bipartite host--MPD topology.
The objective is to minimise the peak per-MPD load over the full trace
window, which directly determines pooling savings:
\begin{equation}
  \text{pooling\_savings}
  = 1 - \frac{\max_j\, c_j^{\text{peak}} \cdot m}{D_{\text{pod}}}
  \label{eq:ps}
\end{equation}
A 1--2\% improvement over greedy is considered strong.

Runs 5--7 revealed three failure modes that motivate this ablation study:

\begin{center}
\small
\renewcommand{\arraystretch}{1.3}
\begin{tabular}{llll}
\toprule
\textbf{Run} & \textbf{Reward} & \textbf{Failure mode} & \textbf{Peak savings} \\
\midrule
6 & ``current'' ($\Delta$peak + variance) &
  Policy collapsed to 1--2 MPDs; entropy $\to 0$ &
  $-0.62$ (harmful) \\
5 & Variant A (local projected peak) &
  Savings plateau; entropy collapse by 2M steps &
  $\sim$1.3\% (flat) \\
7 & Variant A (same, different hyperparams) &
  Same plateau; Q-values rising without improvement &
  $\sim$1.2\% (flat) \\
\bottomrule
\end{tabular}
\end{center}

\noindent
The common thread is \textbf{proxy-objective misalignment}: per-step rewards
optimise a quantity (instantaneous peak change) that diverges from the true
objective (episode-level pooling savings) over 7{,}000+ allocation
decisions.  The Noisy Future Knowledge analysis~\cite{futurek} confirms that
5\% per-step prediction error compounds over 14 days to produce
worse-than-greedy outcomes.

\paragraph{Ablation hypothesis.}
Progressively widening the reward scope---from single MPD, to reachable
neighbourhood, to global pod, to episode-level---will reveal which level of
information is necessary and sufficient.  Each level adds exactly one
dimension of information relative to the previous.

% =========================================================================
\section{Notation}
\label{sec:notation}

\begin{center}
\small
\renewcommand{\arraystretch}{1.2}
\begin{tabular}{ll}
\toprule
\textbf{Symbol} & \textbf{Definition} \\
\midrule
$m$                  & number of MPDs in the pod \\
$N(i)$               & MPDs reachable from host $i$ \\
$H(j)$               & hosts connected to MPD $j$ \\
$c_j(t)$             & current allocated load on MPD $j$ (GB) \\
$D_j(t,W)$           & time-weighted departure relief on MPD $j$ over horizon $W$ \\
$S_j(t)$             & sticky load on MPD $j$ (load expected to remain after departures) \\
$x$                  & arriving VM memory request (GB) \\
$\mathbf{w}$         & allocation weight vector across reachable MPDs \\
$W$                  & lookahead window (timesteps) \\
$D_{\text{pod}}$     & total pod DRAM capacity (GB) \\
$\hat{c}_j(t)$       & projected normalized load before placing the new VM \\
$\hat{c}_j^+(t)$     & projected normalized load after placing the new VM \\
\midrule
$c_{\max}^{\text{post}}$ & $\max_j\, c_j$ after allocation ($\star$) \\
$t^*(s)$             & optimal achievable peak via Dinic solver ($\star$) \\
$\delta_t$           & optimality gap ($\star$) \\
$\Phi(s)$            & potential function for PBRS ($\star$) \\
$\mathrm{PS}$        & episode-level pooling savings ($\star$) \\
\bottomrule
\end{tabular}
\end{center}

\paragraph{Core load projection definitions (self-contained).}
For each MPD $j$, define post-placement load and normalized projected loads as
\begin{align}
  c_j^{\text{post}}(t) &= c_j(t) + x\,w_j, \\
  \hat{c}_j(t) &= \frac{c_j(t) - D_j(t,W)}{D_{\text{pod}}}, \\
  \hat{c}_j^+(t) &= \frac{c_j^{\text{post}}(t) - D_j(t,W)}{D_{\text{pod}}}.
\end{align}

% =========================================================================
\section{MPD Reachability}
\label{sec:reachability}

\subsection{Topology and Reachability}

The pod is a bipartite graph $\mathcal{G} = (\mathcal{H}, \mathcal{M},
\mathcal{E})$ where $\mathcal{H}$ is the set of hosts, $\mathcal{M}$ the
set of MPDs, and $\mathcal{E}$ the set of CXL links.  Host $i$ can
allocate VM memory \emph{only} to its reachable set:
\begin{equation}
  N(i) = \{ j \in \mathcal{M} : (i, j) \in \mathcal{E} \}.
\end{equation}
Symmetrically, MPD $j$ receives allocations only from:
\begin{equation}
  H(j) = \{ i \in \mathcal{H} : (i, j) \in \mathcal{E} \}.
\end{equation}

\paragraph{AG16x6 example.}
The default topology (\texttt{AG16x6\_expander\_quads\_r5\_sym\_fixed})
has 16 hosts and a fan-out of up to $d_{\max} = 8$ MPDs per host.
Not all hosts share the same reachable set, and many MPDs are shared across
multiple hosts ($|H(j)| > 1$).  This asymmetry means a uniform reward
scope is inappropriate: the information available to the agent varies by
host position in the graph.

\subsection{Reachability and the Reward Scope Problem}

At each step $t$, the agent receives host $i_t$'s VM request and selects
weights $\mathbf{w}_t$ over $N(i_t)$.  The action strictly cannot affect
any MPD outside $N(i_t)$.  Yet the true objective,
$\max_j c_j^{\text{peak}}$, is over \emph{all} $m$ MPDs.

This creates a fundamental tension: the reward must evaluate local actions
against a global criterion.  Three design responses are possible:
\begin{enumerate}
  \item \textbf{Ignore the gap} --- reward only the local effect (R1, R2).
        Simple but proxy-misaligned.
  \item \textbf{Add a global proxy term} --- penalise unreachable hot MPDs
        as context (R3).  Widens scope but introduces action-independent
        noise.
  \item \textbf{Evaluate at episode end} --- collapse all local decisions
        into one terminal signal (R4).  Fully aligned but sparse.
\end{enumerate}
The ablation ladder tests whether each response is necessary and
sufficient, or whether the credit-assignment cost of R4 dominates the
alignment benefit.

\subsection{Shared MPDs and Inter-Host Interference}

When $|H(j)| > 1$, MPD $j$ is a shared resource.  If host $i_1$ fills it
near capacity, host $i_2$ (which also has $j \in N(i_2)$) loses effective
allocation headroom at future steps.  A reward that observes only
$N(i_t)$ at step $t$ cannot detect this cross-host interference until it
manifests as a load spike---by which point the damage is done.

The R2 formulation captures the direct post-placement effect on all
$j \in N(i_t)$, including shared MPDs connected to the current host.  What
it misses is the \emph{indirect} effect: that filling a shared MPD
constrains a future host's choices.  R3 and R4 address this at different
levels of fidelity, while the optimal gap (Section~\ref{sec:optgap})
provides a model-free lower bound on how much headroom remains globally.

% =========================================================================
\section{Optimal Gap}
\label{sec:optgap}

The codebase includes an exact per-snapshot solver
(\texttt{octopus/optimal.py}) that computes the minimum achievable peak
MPD load $t^*(s_t)$ for a given vector of host demands and topology.
It uses binary search over a Dinic max-flow feasibility check.
This scalar is a lower bound on any policy's peak load at the current
state---the best any allocator could achieve with full knowledge of
instantaneous demands but no knowledge of future arrivals.

\subsection{Optimality-Gap Penalty}

Define the normalised optimality gap after the current allocation:
\begin{equation}
\boxed{
  \delta_t
  \;=\;
  \frac{c_{\max}^{\text{post}} \;-\; t^*(s_t)}{D_{\text{pod}}}
}
\label{eq:optgap}
\end{equation}

$\delta_t = 0$ means the agent's placement is as good as optimal given
current demands.  $\delta_t > 0$ quantifies wasted peak capacity.
This can be added as a shaping term to any base reward $R_k$:
\begin{equation}
  r_t'
  \;=\;
  R_k(t) \;-\; \beta\,\delta_t,
  \qquad
  \beta \in [0.1,\; 0.5]
  \label{eq:optgap-shaped}
\end{equation}

Unlike R1--R3, this term is grounded in a provably tight lower bound
rather than a heuristic, and it is sensitive to every allocation decision
(not just those that change the global worst MPD).

\subsection{PBRS-Compatible Version}

To guarantee the shaped reward does not change the set of optimal policies
(Section~\ref{sec:pbrs}), define the potential:
\begin{equation}
  \Phi_{\text{opt}}(s)
  \;=\;
  -\,\frac{t^*(s)}{D_{\text{pod}}}
  \label{eq:opt-potential}
\end{equation}

Then $F(s,s') = \gamma\,\Phi_{\text{opt}}(s') - \Phi_{\text{opt}}(s)$
rewards reductions in the optimal lower bound itself.  Combined with the
sparse terminal reward (R4), this is the safest dense signal available:
it pushes the policy toward states where the optimally achievable peak is
low, without encoding assumptions about departures or future arrivals.

\subsection{Implementation Cost}

\texttt{find\_optimal} runs a max-flow per call.  At 1{,}000
steps/episode $\times$ 32 envs, calling it every step would dominate
wall-clock time.  Practical options:
\begin{itemize}
  \item Call every $k = 10$--$50$ steps; hold the last value constant
        between calls (stale-optimal shaping).
  \item Call only at episode end to compute a terminal potential bonus.
  \item Pre-compute $t^*$ offline for a representative grid of demand
        vectors and use nearest-neighbour lookup at training time.
\end{itemize}

% =========================================================================
\section{Literature-Based Reward Modifications}
\label{sec:secondary}

The base rewards R1--R4 (Section~\ref{sec:ablation}) can each be improved
by layering principled modifications drawn from RL theory and systems RL
literature.  This section treats each modification as a composable module
that is independent of the base reward choice.

% -------------------------------------------------------------------------
\subsection{Potential-Based Reward Shaping}
\label{sec:pbrs}

Ng, Harada, and Russell~\cite{ng1999pbrs} proved that an agent's set of
optimal policies is preserved \emph{if and only if} shaping rewards take
the form:
\begin{equation}
  F(s, s')
  \;=\;
  \gamma\,\Phi(s') \;-\; \Phi(s)
  \label{eq:pbrs}
\end{equation}
where $\Phi : \mathcal{S} \to \mathbb{R}$ is an arbitrary potential
function.  Any other dense reward risks inducing a suboptimal policy
relative to the true objective.

R1--R3 are not in this form: they encode transition information that
cannot be expressed as a potential difference.  This is acceptable for an
ablation (we want to test their empirical value), but it means R1--R3 can
in principle game the proxy.  R4 avoids this by using the true objective,
but requires dense shaping to be learnable at $\sim$7{,}000-step horizons.

\paragraph{Candidate potentials for this domain.}
\begin{enumerate}
  \item \textbf{Simple:} $\Phi(s) = -\max_j c_j(t) / D_{\text{pod}}$.
        Cheap (one max over $m$ values per step).  Rewards peak reductions
        in a policy-invariant way.
  \item \textbf{Tight:} $\Phi(s) = -t^*(s) / D_{\text{pod}}$.
        Uses the Dinic solver (Section~\ref{sec:optgap}).  Rewards states
        where the optimal achievable peak is low.  Tighter but more
        expensive.
\end{enumerate}

Both potentials are non-positive and approach zero as the system empties,
so $F(s,s') > 0$ when a placement reduces the relevant peak---the
direction we want.

% -------------------------------------------------------------------------
\subsection{Sub-Episode Decomposition}
\label{sec:subepisodes}

Mao et al.~\cite{mao2017pensieve} (Pensieve) and~\cite{mao2019decima}
(Decima) both demonstrate that sparse task-aligned rewards with shorter
horizons outperform dense proxy rewards with long horizons.  Pensieve
uses per-chunk rewards; Decima uses job-completion rewards rather than
per-step queue-length penalties.

Applied here: the 14-day trace produces episodes of $\sim$7{,}000 steps.
Breaking the trace at fixed 1--2 day intervals (288--576 steps at
5-minute resolution) turns R4 from one terminal signal per 7{,}000 steps
into one signal per 300--600 steps---a 10$\times$ improvement in
credit-assignment density while preserving true-objective alignment.

\paragraph{Implementation.}
At each sub-episode boundary:
\begin{enumerate}
  \item Compute sub-episode pooling savings from the peak MPD load observed
        within that window.
  \item Return the sub-episode savings as a terminal reward.
  \item Reset the peak tracker but \textbf{not} the MPD loads---physical
        state carries over; only the measurement window resets.
\end{enumerate}

% -------------------------------------------------------------------------
\subsection{Counterfactual Difference Rewards}
\label{sec:counterfactual}

Wolpert and Tumer~\cite{wolpert2002diff} introduced difference rewards to
isolate an individual agent's contribution to the global objective by
comparing the actual outcome to a counterfactual baseline action
$a_{\text{base}}$:
\begin{equation}
  D(s, a)
  \;=\;
  G(s^{a}) - G(s^{a_{\text{base}}})
  \label{eq:diff-reward}
\end{equation}
where $G(s) = -\max_j \hat{c}_j(s)$ is the global congestion score and
$a_{\text{base}}$ is a cheap reference (e.g.\ greedy or uniform split).

This is most relevant for R3 (Section~\ref{sec:r3}), where the global
term $\max_{j \notin N(i)} \hat{c}_j$ is action-independent: it penalises
global congestion regardless of what the current agent does.  Replacing it
with a difference reward makes the global term action-dependent and
eliminates the moving-baseline problem that inflates critic variance.

% -------------------------------------------------------------------------
\subsection{Dynamic $\lambda$ Scheduling}
\label{sec:lambda}

R3 introduces $\lambda$ to weight the global term against the local term.
A fixed $\lambda$ is problematic: the right trade-off changes during
training, and the topology does not support uniform load.  Three
alternatives:

\begin{center}
\small
\renewcommand{\arraystretch}{1.3}
\begin{tabular}{p{3.5cm}p{4.5cm}p{4.5cm}}
\toprule
\textbf{Approach} & \textbf{Expected behaviour} & \textbf{Complexity} \\
\midrule
Remove entirely (use R2) &
  Eliminates critic noise from action-independent global term &
  Simplest \\
Dynamic scheduling &
  Start $\lambda = 0$ (pure local); increase as policy stabilises
  via bandit over validation pooling savings &
  Moderate; requires scheduler \\
Lexicographic &
  Optimise local peak first; use global only as tiebreaker when
  local peaks are equal &
  Moderate; two-head critic or constraint \\
\bottomrule
\end{tabular}
\end{center}

\noindent
For the ablation, run R3 both with and without $\lambda$ to isolate its
contribution.  If R2 + PBRS + R4 terminal matches or exceeds R3, then
$\lambda$ adds complexity without value and should be dropped.

% =========================================================================
\section{Ablation Formulations}
\label{sec:ablation}

Six formulations, each adding exactly one dimension of information or one
design correction over the previous.  R1--R4 form a clean scope
progression from narrowest signal to true objective.  The final two
variants layer the modifications from Section~\ref{sec:secondary} onto
the objective-aligned base.

% -------------------------------------------------------------------------
\subsection{R1: Single MPD Reward}
\label{sec:r1}

\paragraph{Definition.}
Penalise only the single worst MPD after placement.
\begin{equation}
\boxed{
  R_1(t)
  \;=\;
  -\,\frac{c_{j^*}^{\text{post}}(t)}{D_{\text{pod}}},
  \qquad
  j^* = \arg\max_{j}\, c_j^{\text{post}}(t)
}
\label{eq:r1}
\end{equation}

\paragraph{Rationale.}
The simplest possible signal: how bad is the single worst MPD right now?
No neighbourhood context, no departure awareness, no global view.  Serves
as the controlled floor of the ablation---if any other variant cannot beat
R1, it is adding noise rather than signal.

Unlike Run~6's ``current'' reward, R1 omits the variance term and
fair-share normalisation, isolating the effect of scope (single vs.\
group) from the effect of mixing incommensurable objectives (peak vs.\
variance).

\paragraph{Signal properties.}
\begin{itemize}
  \item \textbf{Density:} Dense. \quad \textbf{Scope:} Single global worst
        MPD. \quad \textbf{Range:} $(-1,\; 0]$. \quad
        \textbf{Departure-aware:} No.
\end{itemize}

\paragraph{Expected failure mode.}
When the allocation does not land on the current worst MPD \emph{and} does
not create a new worst, the reward is identical regardless of which MPD is
chosen.  Policy gradient is zero with respect to action choice in these
cases---the same mechanism that caused Run~6's entropy collapse.

% -------------------------------------------------------------------------
\subsection{R2: Localized Group Reward}
\label{sec:r2}

\paragraph{Definition.}
Penalise the worst projected load across all MPDs reachable from the
current host.
\begin{equation}
\boxed{
  R_2(t)
  \;=\;
  -\,\max_{j \,\in\, N(i)}\, \hat{c}_j^{\,+}(t)
}
\label{eq:r2}
\end{equation}

\paragraph{Rationale.}
R2 improves on R1 in two ways.  First, all reachable MPDs contribute: the
agent receives distinct signals for placing on a lightly-loaded vs.\
heavily-loaded reachable MPD, even if neither is the global worst.  Second,
the projected load subtracts departure relief $D_j$, distinguishing ``full
but draining soon'' from ``full and staying full.''  This departure
awareness is the most important design improvement over R1.

\paragraph{Signal properties.}
\begin{itemize}
  \item \textbf{Density:} Dense. \quad \textbf{Scope:} All $j \in N(i)$
        (typically $d_{\max} = 8$). \quad \textbf{Range:} $(-1,\; 0]$.
        \quad \textbf{Departure-aware:} Yes.
\end{itemize}

\paragraph{Expected failure mode.}
R2 is blind to unreachable MPDs.  A neighbourhood where all MPDs are at
60\% looks the same as one where one MPD is at 60\% and the rest are
empty---R2 loses distributional information.  Cross-host interference on
shared MPDs (Section~\ref{sec:reachability}) is also invisible until it
manifests as a spike.

% -------------------------------------------------------------------------
\subsection{R3: Global Group Reward}
\label{sec:r3}

\paragraph{Definition.}
Penalise the worst projected load locally, plus a weighted penalty for
the worst unreachable MPD.
\begin{equation}
\boxed{
  R_3(t)
  \;=\;
  -\,\max_{j \,\in\, N(i)}\, \hat{c}_j^{\,+}(t)
  \;-\;
  \lambda \cdot \max_{j \,\notin\, N(i)}\, \hat{c}_j(t)
}
\label{eq:r3}
\end{equation}

\paragraph{Rationale.}
R3 tests whether the agent benefits from seeing global pod state beyond
its neighbourhood.  The second term penalises hot unreachable MPDs,
signalling that neighbouring hosts are under stress.  Even though the
current host cannot directly relieve that stress, the value function may
learn to associate global stress with worse long-term outcomes.

\paragraph{Action-dependent variants.}
The baseline R3 (Eq.~\ref{eq:r3}) suffers from an action-independent
global term.  Applying the Wolpert--Tumer difference reward
(Section~\ref{sec:counterfactual}) yields action-dependent alternatives:
\begin{align}
\textbf{R3-C:}\quad
R_{3\text{C}}(t)
&=
-\,\max_{j \in N(i)} \hat{c}_j^{+}(t)
- \lambda \bigl(G(s_{t+1}) - G(s_t)\bigr),
\label{eq:r3c}
\\[4pt]
\textbf{R3-D:}\quad
R_{3\text{D}}(t)
&=
R_2(t)
+ \lambda \bigl(G(s_{t+1}^{a_t}) - G(s_{t+1}^{a_{\text{base}}})\bigr),
\label{eq:r3d}
\end{align}
where $G(s) = -\max_j \hat{c}_j(s)$ and $a_{\text{base}}$ is a greedy or
uniform-split baseline.  R3-D provides the highest-quality credit
assignment for global effect but requires an extra rollout per step.

\paragraph{Signal properties.}
\begin{itemize}
  \item \textbf{Density:} Dense. \quad \textbf{Scope:} All $m$ MPDs.
        \quad \textbf{Range:} $(-(1{+}\lambda),\; 0]$. \quad
        \textbf{Departure-aware:} Yes (local term and $\hat{c}_j$ for
        unreachable MPDs).
\end{itemize}

\paragraph{Expected failure mode.}
R3 and R3-C still suffer from partial action independence; the global
component can become a moving baseline that increases critic variance
without improving action ranking.  R3-D fixes this in principle but
introduces non-trivial compute overhead and reward-scale sensitivity in
the choice of $\lambda$.

% -------------------------------------------------------------------------
\subsection{R4: Episode-Level Reward}
\label{sec:r4}

\paragraph{Definition.}
Zero intermediate reward; true pooling savings at episode end.
\begin{equation}
\boxed{
  R_4(t)
  \;=\;
  \begin{cases}
    0 & t < T \\[4pt]
    \mathrm{PS}(\text{episode}) & t = T
  \end{cases}
}
\label{eq:r4}
\end{equation}

\paragraph{Rationale.}
R4 is the only base formulation that directly optimises the true
objective.  No proxy, no $\lambda$, no intermediate approximation.
Every distortion introduced by R1--R3 is eliminated.  The agent receives
a single scalar at episode end that is exactly the quantity we want to
maximise.

\paragraph{Signal properties.}
\begin{itemize}
  \item \textbf{Density:} Sparse ($\sim$1 signal per 7{,}000 steps).
        \quad \textbf{Scope:} Global and temporal (captures all placements
        over the full trace). \quad \textbf{Range:} $(-1,\; 1)$ in
        practice. \quad \textbf{Departure-aware:} Not directly (terminal
        reward reflects actual peaks, not predicted ones).
\end{itemize}

\paragraph{Expected failure mode.}
Credit assignment.  With 7{,}000 steps and one terminal reward, the critic
must learn which of the 7{,}000 decisions mattered.  SAC's replay buffer
and Bellman bootstrapping can propagate this signal backward, but
convergence is slow without dense shaping.  This motivates the next two
variants.

\paragraph{Implementation note.}
Episode-level pooling savings is already computed in the codebase
(\texttt{pooling\_ratio} / \texttt{pooling\_savings} in the eval callback).
It needs to be wired as the terminal reward rather than a separate logging
metric.

% -------------------------------------------------------------------------
\subsection{R4 + PBRS + Sub-Episodes}
\label{sec:r4pbrs}

\paragraph{Definition.}
Episode-level terminal reward with dense potential-based shaping and a
shorter sub-episode horizon for credit assignment.
\begin{equation}
\boxed{
  r_t
  \;=\;
  \begin{cases}
    \gamma\,\Phi(s_{t+1}) - \Phi(s_t)
      & t \text{ not a sub-episode boundary} \\[6pt]
    \mathrm{PS}_{\text{sub}}(t)
      + \gamma\,\Phi(s_{t+1}) - \Phi(s_t)
      & t \text{ a sub-episode boundary}
  \end{cases}
}
\label{eq:r4pbrs}
\end{equation}
where $\Phi(s) = -\max_j c_j / D_{\text{pod}}$ (simple potential) or
$\Phi(s) = -t^*(s)/D_{\text{pod}}$ (tight potential), and
$\mathrm{PS}_{\text{sub}}(t)$ is the pooling savings over the current
sub-episode window (1--2 days $\approx$ 288--576 steps).

\paragraph{Rationale.}
This variant stacks the two literature-based improvements from
Sections~\ref{sec:pbrs} and~\ref{sec:subepisodes}.  PBRS (Ng et al.)
provides dense per-step signal without distorting the optimal policy.
Sub-episode decomposition (Pensieve/Decima) shortens the credit-assignment
horizon from $\sim$7{,}000 steps to $\sim$300--600 steps.  Together they
address the sole failure mode of R4 (sparse signal) while preserving full
objective alignment.

\paragraph{Signal properties.}
\begin{itemize}
  \item \textbf{Density:} Dense (per-step PBRS shaping) plus terminal at
        each sub-episode boundary.
  \item \textbf{Scope:} Global and temporal (terminal captures all
        placements in each sub-episode window).
  \item \textbf{Objective-aligned:} Yes---PBRS preserves the optimal policy
        set; sub-episode savings is a true partial objective.
\end{itemize}

\paragraph{Expected failure mode.}
Added implementation complexity at sub-episode boundaries (peak tracker
reset without MPD state reset).  PBRS and sub-episodes may be partially
redundant: if PBRS alone is sufficient, sub-episodes add complexity without
benefit.  The ablation should test R4+PBRS alone and the full stack
separately to separate their contributions.

% -------------------------------------------------------------------------
\subsection{Best Variant + Optimal Gap Shaping}
\label{sec:r4opt}

\paragraph{Definition.}
The best-performing variant from the preceding five, augmented with the
optimality-gap penalty (Section~\ref{sec:optgap}).  Plain penalty form:
\begin{equation}
\boxed{
  r_t^+
  \;=\;
  r_t^{\text{best}} \;-\; \beta\,\delta_t,
  \qquad
  \delta_t = \frac{c_{\max}^{\text{post}} - t^*(s_t)}{D_{\text{pod}}},
  \quad \beta \in [0.1,\; 0.5]
}
\label{eq:r4opt}
\end{equation}

PBRS-compatible form using the tight potential:
\begin{equation}
  r_t^+
  \;=\;
  r_t^{\text{best}}
  \;+\;
  \gamma\,\Phi_{\text{opt}}(s_{t+1}) - \Phi_{\text{opt}}(s_t),
  \qquad
  \Phi_{\text{opt}}(s) = -\,\frac{t^*(s)}{D_{\text{pod}}}
  \label{eq:r4opt-pbrs}
\end{equation}

\paragraph{Rationale.}
The optimal gap term is the only shaping signal grounded in a provably
tight lower bound (Dinic max-flow).  It rewards the agent for placing VMs
close to the theoretical optimum given current state, independent of
future arrivals.  Applied to whichever base variant already achieves the
best objective alignment, it provides an additional per-step calibration
signal at the cost of one max-flow solve per call.

The benefit is largest on R1 and R2 (weakest signal; $\delta_t$ adds a
lower-bound reference they completely lack) and smallest on R4 and R4+PBRS
(already optimising the true objective).

\paragraph{When to evaluate.}
Optimal-gap shaping should be the last variant tested, after the
five-variant primary sweep is complete.  The computational overhead is
non-trivial and only justified once the strongest objective-aligned base
is identified.

\paragraph{Signal properties.}
\begin{itemize}
  \item \textbf{Density:} Dense (or stale-optimal at every $k$ steps).
        \quad \textbf{Scope:} Global (Dinic solver sees full topology).
        \quad \textbf{Objective-aligned:} Yes (PBRS-compatible form).
\end{itemize}

% =========================================================================
\section{Experimental Plan}
\label{sec:experiment}

\subsection{Ablation Summary}

\begin{center}
\small
\renewcommand{\arraystretch}{1.3}
\begin{tabular}{p{3.2cm}p{2.8cm}ll p{3.8cm}}
\toprule
\textbf{Variant} & \textbf{Scope} & \textbf{Dense?} &
  \textbf{Objective-aligned?} & \textbf{Key risk} \\
\midrule
R1 & Single worst MPD & Yes & No &
  Zero gradient for non-worst placements \\
R2 & Reachable-MPD group & Yes & No &
  Blind to shared-MPD congestion \\
R3 & Full pod (proxy) & Yes & No &
  Action-independent global term \\
R4 & Full episode & No & Yes &
  7{,}000-step credit assignment \\
R4 + PBRS + Sub & Full episode & Yes & Yes &
  Boundary complexity; possible PBRS/sub redundancy \\
Best + Opt.\ gap & Full pod + tight lb & Yes & Yes &
  Compute cost of Dinic per step \\
\bottomrule
\end{tabular}
\end{center}

\subsection{Run Order}

\begin{enumerate}
  \item \textbf{R1} --- Establish the controlled floor.  Cheapest run and
        cleanest negative control.
  \item \textbf{R2} --- Measure how much local projected-load awareness
        buys over R1.
  \item \textbf{R3} --- Test whether global context (action-independent)
        adds value over R2.
  \item \textbf{R4} --- Test whether sparse true-objective signal can
        learn at all without shaping.
  \item \textbf{R4 + PBRS + Sub-Episodes} --- Test whether objective
        alignment plus improved credit assignment is the winning
        combination.
  \item \textbf{Best + Optimal Gap} --- Only after the winner of the
        above five is clearly established; adds non-trivial compute.
\end{enumerate}

\subsection{Metrics}

\begin{itemize}
  \item \textbf{Primary:} \texttt{eval/pooling\_savings\_mean} on held-out
        trace (10 random pod mappings).
  \item \textbf{Diagnostic:} entropy coefficient
        (\texttt{sac/ent\_coef}), critic loss convergence speed, Q-value
        growth rate.
  \item \textbf{Cost:} wall-clock time per 1M steps.
\end{itemize}

\subsection{Success Criteria}

A variant ``beats'' another if it achieves $> 0.5\%$ higher pooling
savings on the held-out eval trace, sustained for at least 3 consecutive
eval checkpoints (300k steps).  Ties are broken by lower variance across
pod mappings.

% =========================================================================
\section{Empirical Results: Reward Ablation on Held-Out Traces}
\label{sec:results}

All variants trained for 500k steps with SAC, $n_\text{envs}=32$, multi-trace
rotation, augmentation enabled, HOTFIX filter off.
Evaluation: 20 pod seeds per trace, topology \texttt{AG16x6\_expander\_quads\_r5\_sym\_fixed}.

\subsection{LVL01PrdApp05 (churn-dominated trace)}

Greedy allocation is harmful on this trace: pooling ratio $>1$ means
the policy must allocate CXL memory that exceeds the savings from pooling.
All RL variants improve over greedy.

\begin{table}[H]
\centering
\begin{tabular}{lrrr}
\toprule
\textbf{Policy} & \textbf{Pool ratio} & \textbf{Savings} & \textbf{vs.\ Greedy} \\
\midrule
Greedy          & 1.033 & $-3.30\%$ & --- \\
\textbf{R2}     & \textbf{1.002} & $\mathbf{-0.20\%}$ & $\mathbf{+3.1\,\text{pp}}$ \\
R1              & 1.008 & $-0.76\%$ & $+2.5\,\text{pp}$ \\
R4              & 1.009 & $-0.88\%$ & $+2.4\,\text{pp}$ \\
R3              & 1.010 & $-0.96\%$ & $+2.3\,\text{pp}$ \\
current         & 1.013 & $-1.31\%$ & $+2.0\,\text{pp}$ \\
R5              & 1.014 & $-1.40\%$ & $+1.9\,\text{pp}$ \\
\bottomrule
\end{tabular}
\caption{LVL01 held-out evaluation (20 seeds, $n=20$ iterations).
  Lower pool ratio is better. All RL variants beat greedy; R2 is closest to break-even.}
\label{tab:lvl01}
\end{table}

\subsection{AMS20PrdApp19 (long-lived VM trace)}

Greedy performs strongly on this trace (ratio $0.721$, savings $+27.9\%$).
All RL variants cluster near $6.5\%$ savings — a $\approx21\,\text{pp}$ gap behind greedy.
Reward variant makes essentially no difference.

\begin{table}[H]
\centering
\begin{tabular}{lrrr}
\toprule
\textbf{Policy} & \textbf{Pool ratio} & \textbf{Savings} & \textbf{vs.\ Greedy} \\
\midrule
Greedy          & 0.721 & $+27.9\%$ & --- \\
R1              & 0.934 & $+6.6\%$  & $-21.2\,\text{pp}$ \\
current         & 0.934 & $+6.6\%$  & $-21.2\,\text{pp}$ \\
R2              & 0.935 & $+6.6\%$  & $-21.3\,\text{pp}$ \\
R4              & 0.934 & $+6.6\%$  & $-21.3\,\text{pp}$ \\
R3              & 0.936 & $+6.4\%$  & $-21.5\,\text{pp}$ \\
R5              & 0.937 & $+6.3\%$  & $-21.6\,\text{pp}$ \\
\bottomrule
\end{tabular}
\caption{AMS20 held-out evaluation (20 seeds, $n=20$ iterations).
  All RL variants are statistically indistinguishable from each other
  and trail greedy by $\approx21\,\text{pp}$. Reward shaping has no effect on this trace.}
\label{tab:ams20}
\end{table}

\subsection{Cross-trace summary}

The two held-out traces reveal a generalisation gap rather than a reward-shaping
effect. On LVL01 (high churn: 64\% of VMs live $<1$\,h, median lifetime
$\approx$30\,min) every RL variant outperforms greedy. On AMS20 (long-lived:
46\% of VMs live $>7$\,days, mean $\approx$3.8\,days) greedy dominates.
Reward variant R1--R5 differences are within noise on both traces,
suggesting the bottleneck is \emph{generalisation across VM-lifetime
distributions}, not the reward formulation.

% =========================================================================
\section{Trace Analysis}
\label{sec:trace-analysis}

\subsection{Ghost VMs (data quality)}

Both traces contain VMs whose \texttt{start\_time} and \texttt{end\_time}
fields were never populated during trace loading (they retain the class
defaults: \texttt{start=2024-01-01}, \texttt{end=2020-01-01}).
AMS20 has 4{,}788 such ghost VMs (14.1\%); LVL01 has 50{,}674 (18.7\%).

These are \textbf{safely filtered} by \texttt{\_build\_events}: the base
timestamp for each episode is derived from the real VM population
($\approx2023$), so the ghost end-time (2020) maps to a tick far
below zero, triggering the \texttt{if vm\_e < 0: continue} guard before
any event is appended.  Ghost VMs generate no events and do not corrupt
the simulation.

\subsection{VM-lifetime distributions across all ten traces}

Table~\ref{tab:traces} characterises the full training-trace pool.
AMS20 is the only predominantly long-lived trace: 46\% of valid VMs
run for more than seven days and the median lifetime is $\approx$20\,h.
The remaining nine traces are churn-dominated, with LVL01 being the
most extreme (64\% of VMs live less than one hour, median $\approx$30\,min).

\begin{table}[H]
\centering
\small
\begin{tabular}{lrrrrrrrr}
\toprule
\textbf{Trace} & \textbf{Valid VMs} & \textbf{Ghost\%} & \textbf{$<$1h\%} & \textbf{$>$7d\%} & \textbf{Med (h)} & \textbf{Mean (d)} & \textbf{AvgConc} & \textbf{Med GB} \\
\midrule
\textbf{AMS20}   & 27{,}159  & 20.2 & 22.3 & \textbf{49.3} & \textbf{124.6} & \textbf{4.14} & 14{,}053 & 8  \\
BLA              & 100{,}404 & 22.6 & 29.8 & 31.2 &   6.2 & 2.79 & 35{,}032 & 8  \\
BN9              & 138{,}335 & 22.4 & 30.5 & 21.1 &   2.8 & 1.92 & 33{,}231 & 16 \\
DSM08            &  83{,}252 & 19.8 & 33.8 & 26.3 &   2.8 & 2.56 & 26{,}677 & 7  \\
DUB24            &  64{,}493 & 18.5 & 24.9 & 37.5 &  22.8 & 3.43 & 27{,}629 & 7  \\
LON23            &  17{,}831 & 24.0 & 35.1 & 28.3 &   2.8 & 2.47 &  5{,}498 & 14 \\
\textbf{LVL01}   & 181{,}558 & 33.0 & \textbf{56.7} & 10.5 & \textbf{0.75} & \textbf{0.98} & 22{,}334 & 14 \\
SG2              & 109{,}483 & 22.4 & 32.7 & 29.2 &   2.3 & 2.52 & 34{,}498 & 14 \\
SYD21            &  46{,}496 & 19.1 & 26.6 & 36.5 &   7.2 & 3.13 & 18{,}181 & 16 \\
YTO21            &  58{,}200 & 27.5 & 36.4 & 30.5 &   1.3 & 2.56 & 18{,}606 & 8  \\
\bottomrule
\end{tabular}
\caption{VM-lifetime statistics across the ten training traces.
  Ghost\% = fraction with unset timestamps (safely filtered before simulation).
  AMS20 is a clear outlier: median lifetime 5+ days, 49\% of VMs survive beyond the 7-day window.
  LVL01 is the most churn-heavy: median 45 min, 57\% of VMs live under 1 hour.
  The 9-to-1 skew toward churn-dominated traces in the training pool
  explains the generalisation gap observed on AMS20.}
\label{tab:traces}
\end{table}

\subsection{Training composition and the generalisation gap}

Both AMS20 and LVL01 are included in the ten-trace multi-trace training
pool (\texttt{--multi-trace}).  Despite this, the RL policy achieves only
$\approx6.5\%$ savings on AMS20 versus greedy's $27.9\%$.
The gap is not caused by AMS20 being unseen during training.

The structural explanation is that on long-lived traces the greedy
heuristic (always assign to the least-loaded MPD) is already near-optimal:
once the initial allocation is made it persists for days, leaving little
room for a learned policy to improve.  On churn-dominated traces, rapid
reallocation opportunities arise that greedy cannot exploit, giving the RL
policy its advantage.

A lifetime-conditioned oracle baseline (assign short-lived VMs greedily,
use optimal placement only for long-lived VMs) was evaluated to quantify
how much of the AMS20 gap is structurally recoverable.

\subsection{Lifetime-conditioned baseline}

The baseline routes each VM arrival based on its residual lifetime:
VMs with $\text{dealloc\_tick} - \text{tick} \geq 288$ (24\,h) are
allocated to the \emph{most-loaded} accessible MPD (bin-pack / consolidate);
shorter-lived VMs use the standard greedy water-fill.

\begin{table}[H]
\centering
\begin{tabular}{lrrr}
\toprule
\textbf{Trace} & \textbf{Greedy savings} & \textbf{Lifetime-cond.} & \textbf{$\Delta$} \\
\midrule
AMS20 & $+28.6\%$ & $-567\%$ & $-596\,\text{pp}$ \\
LVL01 & $-3.3\%$  & $-664\%$ & $-661\,\text{pp}$ \\
\bottomrule
\end{tabular}
\caption{Lifetime-conditioned baseline (24\,h threshold, 20 seeds).
  Anti-greedy allocation for long-lived VMs is catastrophically worse on both traces.}
\label{tab:lifetime}
\end{table}

The result confirms a key invariant of the pooling objective: greedy
water-fill \emph{minimises the instantaneous peak MPD load} for any single
VM arrival, regardless of the VM's lifetime.  Anti-greedy (consolidating
to the most-loaded MPD) concentrates load and maximises the peak.
No causal strategy that deviates from greedy on individual allocations
can improve on greedy's instantaneous optimality.

\paragraph{Implication.}
The 21\,pp gap between RL and greedy on AMS20 is \emph{structurally
unrecoverable} by any online policy.  For long-lived workloads, each
placement decision persists for days and greedy's myopic optimality is
also globally optimal---there are no short-term deviations that pay off
long-term.  The RL agents underperform greedy on AMS20 because they
have internalised churn-oriented heuristics (learned from the 9 churn-heavy
training traces) that actively hurt when VMs are long-lived.
Reward shaping (R1--R5) cannot fix this; the bottleneck is the
training distribution, not the reward signal.

% =========================================================================
\begin{thebibliography}{9}

\bibitem{ng1999pbrs}
A.~Ng, D.~Harada, S.~Russell.
\newblock Policy Invariance Under Reward Transformations: Theory and Application
  to Reward Shaping.
\newblock In \textit{ICML}, 1999.

\bibitem{pond2023}
Huaicheng Li, Daniel S.\ Berger, Lisa Hsu, et al.
\newblock Pond: {CXL}-Based Memory Pooling Systems for Cloud Platforms.
\newblock In \textit{ASPLOS}, 2023.

\bibitem{futurek}
Richard Uzeel.
\newblock Perfect Future Knowledge.
\newblock Internal project report, 2026.

\bibitem{mao2017pensieve}
H.~Mao et al.
\newblock Real World Performance of Adaptive Bitrate Algorithms (Pensieve).
\newblock In \textit{ACM SIGCOMM}, 2017.

\bibitem{mao2019decima}
H.~Mao et al.
\newblock Learning Scheduling Algorithms for Data Processing Clusters (Decima).
\newblock In \textit{ACM SIGCOMM}, 2019.

\bibitem{wolpert2002diff}
D.~H.~Wolpert, K.~Tumer.
\newblock Collective Intelligence, Data Routing and Braess' Paradox.
\newblock \textit{Journal of Artificial Intelligence Research}, 16:359--387, 2002.

\bibitem{pbrs2025}
Improving the Effectiveness of Potential-Based Reward Shaping in Reinforcement
  Learning.
\newblock arXiv, 2025.

\end{thebibliography}

\end{document}
